\subsection{Układ pracy}\label{section_layout}
\quad Praca została podzielona na [X] rozdziałów.  W pierwszym przedstawiono problematykę badawczą oraz cele pracy, a także dokonano przeglądu literatury w zakresie automatyzacji planowania zadań oraz modelowania energii człowieka i rytmu dobowego. Drugi rozdział przedstawia teoretyczne wprowadzenie do rozważanego problemu oraz omówienie wybranych metod. W następnych rozdziałach omówiono kolejno strukturę danych oraz konstrukcję modeli i przeprowadzanego eksperymentu. Ostatecznie w rozdziałach \ref{section_results} i \ref{section_summary} przedstawiono i omówiono wyniki, przeprowadzając ocenę skuteczności wybranych metod, a także uwzględniono możliwe rozwinięcie tematyki personalizowanego planowania zadań. 
% Praca została podzielona na siedem rozdziałów. W pierwszym przedstawiono problematykę badawczą oraz cele pracy, a także dokonano przeglądu literatury dotyczącej systemów rekomendacyjnych. Drugi rozdział zawiera opis wybranych algorytmów oraz podstawowe zagadnienia związane z rekomendacjami. W kolejnym rozdziale przeanalizowano i porównano dostępne aplikacje służące do rekomendacji kierunków studiów. Rozdział czwarty opisuje technologie wykorzystane przy implementacji opracowanego systemu, natomiast piąty szczegółowo omawia proces jego tworzenia, poszczególne komponenty oraz funkcjonalności. Szósty rozdział pełni rolę podręcznika użytkownika, zawierając dokładny opis interfejsu aplikacji, a siódmy stanowi podsumowanie projektu, omawia uzyskane rezultaty oraz wskazuje możliwości dalszego rozwoju. 